\documentclass[12pt,a4paper]{article}

% Packages essentiels
\usepackage[utf8]{inputenc}
\usepackage[T1]{fontenc}
\usepackage[french]{babel}
\usepackage{amsmath}
\usepackage{amssymb}
\usepackage{geometry}
\usepackage{hyperref}

% Configuration de la page
\geometry{a4paper, margin=2.5cm}
\hypersetup{
    colorlinks=true,
    linkcolor=blue,
    filecolor=magenta,      
    urlcolor=cyan,
}

% En-tête du document
\title{\textbf{Plan de Recherche :}\\ Étude des matériaux ferroélectriques, composites PZT/PVDF et propagation de fissures}
\author{Maxime Defaux \\ ENSEIRB-MATMECA}
\date{\today}

\begin{document}

\maketitle

\tableofcontents
\newpage

\section{Étude des matériaux ferroélectriques (Fondements théoriques)}

Cette section établit le cadre thermodynamique et cristallographique rigoureux à l'échelle du monocristal et du polycristal homogène.

\subsection{Origine microscopique et physique du solide}
\begin{itemize}
    \item Structure cristalline des pérovskites (cas du $\text{BaTiO}_3$ ou du PZT en phase tétragonale/rhomboédrique).
    \item Brisure de symétrie, déplacement ionique et apparition du dipôle permanent.
    \item Notion de domaines ferroélectriques (parois à $90^\circ$ et $180^\circ$) et minimisation de l'énergie de gradient.
\end{itemize}

\subsection{Thermodynamique des transitions de phase (Modèle de Landau-Devonshire)}
\begin{itemize}
    \item Expression de l'énergie libre de Gibbs locale en 3D en fonction du vecteur polarisation $P_i$ et du tenseur des contraintes $\sigma_{ij}$. Le développement polynomial s'exprime sous la forme :
    $$G(P, T) = G_0 + \frac{1}{2}\alpha(T - T_c)P^2 + \frac{1}{4}\beta P^4 + \frac{1}{6}\gamma P^6 - EP$$
    \item Intégration du couplage électromécanique d'ordre supérieur (électrostriction $Q_{ijkl}P_iP_j\sigma_{kl}$).
    \item Analyse des puits de potentiel et prédiction de la polarisation spontanée $P_s$.
\end{itemize}

\subsection{Lois de comportement macroscopiques et non-linéarités}
\begin{itemize}
    \item Équations constitutives du comportement piézoélectrique linéaire :
    $$\sigma_{ij} = C_{ijkl} \varepsilon_{kl} - e_{kij} E_k$$
    $$D_i = e_{ikl} \varepsilon_{kl} + \epsilon_{ij}^S E_j$$
    \item Modélisation de la saturation, du cycle d'hystérésis et de la déformation en papillon.
    \item Formulation des critères de basculement de domaine via la thermodynamique des processus irréversibles.
\end{itemize}

\section{Composites ferroélectriques (PZT + PVDF) : Approche Multi-échelle}

Étude de l'association d'une céramique ferroélectrique rigide (PZT) avec un polymère ferroélectrique souple (PVDF) en connectivité 0-3.

\subsection{Propriétés des phases en présence et connectivité}
\begin{itemize}
    \item Propriétés intrinsèques du PZT (rigidité, $d_{33}$ élevé, forte permittivité).
    \item Propriétés intrinsèques du PVDF (faible rigidité, grande déformabilité, faible permittivité).
\end{itemize}

\subsection{Homogénéisation électromécanique (Micro-Mécanique)}
\begin{itemize}
    \item Formulation du problème de localisation pour les milieux couplés piézoélectriques.
    \item Application des modèles analytiques (Mori-Tanaka adapté, inclusion d'Eshelby couplée).
    \item Détermination des tenseurs effectifs $C_{ijkl}^*$, $e_{kij}^*$ et $\epsilon_{ij}^*$ en fonction de la fraction volumique de PZT.
\end{itemize}

\subsection{Problématique de la polarisation du composite (Polling)}
\begin{itemize}
    \item Analyse du saut de permittivité entre les deux phases.
    \item Modélisation du phénomène de blindage du champ électrique au sein de la matrice PVDF.
\end{itemize}

\section{Propagation des fissures (Mécanique de la rupture couplée)}

Analyse de la durabilité et de la rupture exacerbée par les concentrations de champs électromécaniques.

\subsection{Mécanique de la rupture en milieu piézoélectrique linéaire}
\begin{itemize}
    \item Définition des facteurs d'intensité de contraintes ($K_I, K_{II}$) et de déplacement électrique ($K_D$).
    \item Calcul du taux de restitution d'énergie via l'extension de l'intégrale $J$ aux milieux couplés :
    $$J = \int_{\Gamma} \left( W n_1 - t_i \frac{\partial u_i}{\partial x_1} - q \frac{\partial \phi}{\partial x_1} \right) d\Gamma$$
\end{itemize}

\subsection{Effet du champ électrique sur la ténacité}
\begin{itemize}
    \item Asymétrie de la rupture selon l'orientation du champ appliqué.
    \item Explication mécanique de l'ouverture ou de la fermeture induite par l'effet piézoélectrique.
\end{itemize}

\subsection{Non-linéarités en pointe de fissure : Zone de basculement de domaine}
\begin{itemize}
    \item Modélisation de la zone de concentration (\textit{process zone}) provoquant un basculement local.
    \item Effet de blindage ou d'amplification de la fissure (plasticité de transformation).
    \item Stratégie de modélisation numérique (XFEM ou \textit{Phase-Field}) pour le suivi de fissure dans la microstructure.
\end{itemize}

\end{document}